\section{Probabilidades}

\subsection{A Teoria de Probabilidades}
Ao observar, registrar ou extrair uma amostra de dados a partir de um fenômeno aleatório,
isto é, cujo resultado não pode ser deterministicamente previsto, a Estatística, dentre
outras coisas, procura entender o comportamento das frequências de ocorrência dos valores
amostrados.
Essas frequências são estimativas das \textbf{probabilidades} dos eventos observados.
A partir de hipóteses adequadas, podemos modelar as frequências por modelos probabilísticos
que, por sua vez, são objetos de estudo da Teoria das Probabilidades.

A Teoria de Probabilidades encontra seus fundamentos em alguns campos da matemática pura,
tais como a Matemática Discreta, Análise e Teoria da Medida.
No entanto, neste texto, estaremos mais interessados na visão da teoria para aplicações
estatísticas, como a já a modelagem de frequências e estimação de parâmetros de modelos.

Probabilidades são frequentemente representadas como porcentagens e são muitas vezes
entendidas como a frequência de um evento quando se repete o experimento um número
suficiente de vezes. Por exemplo, dizer que 16\% da população possui diploma de
graduação significaria que se escolhermos aleatoriamente 100 pessoas dessa população,
cerca de 16 delas teriam diploma teriam diploma de graduação. Essa abordagem é chamada
de \textbf{frequentista}, sendo a mais amplamente adotada em aplicações estatísticas
clássicas.

Outra abordagem para compreender probabilidades é enquanto o grau de crença numa hipótese,
como ao dizer que se tem 80\% de certeza em dizer que irá chover, visto as evidências do
céu nublado e ventos frios. Essa é uma visão mais intuitiva e moderna de probabilidades,
sendo chamada de \textbf{Bayesiana}, sendo a base da Inferência Bayesiana.

\subsection{$\sigma$-Álgebras}

A noção primitiva que iremos usar como base será de experimento ou fenômeno aleatório,
que modela qualquer fenômeno cujos resultados são conhecidos, porém podem ser
distintos apesar da repetição de condições iniciais.

\begin{def:espaco amostral}
O espaço amostral $\Omega$ de um experimento aleatório é o conjunto de todos os resultados
possíveis do experimento.
\end{def:espaco amostral}

Alguns subconjuntos de um espaço amostral compõe um \textbf{evento} do experimento
aleatório.
Intuitivamente, para serem ``bem-comportados'', os eventos devem satisfazer
algumas regras que possibilitem determinar outros eventos ``bem comportados''.
Formalmente, os eventos devem compor uma $\sigma$-Álgebra.

\begin{def:sigma algebra}
Seja $\Omega$ um conjunto e $\mathcal{B} \subset \Parts{\Omega}$.
O conjunto $\mathcal{B}$ é uma $\sigma$-Álgebra de $\Omega$ se, e somente se,
\begin{enumerate}
	\item[i] $\varnothing \in \mathcal{B}$.
	\item[ii] $\forall A \in \mathcal{B}, A^c \in \mathcal{B}$.
	\item[iii] $\forall A \in \mathcal{B},\ \bigcup_{i=1}^\infty A_i \in \mathcal{B}$.
\end{enumerate}
Se $\Omega$ é um espaço amostral, então $A \in \mathcal{B}$ é denominado evento de
$\Omega$.
\end{def:sigma algebra}

Em resumo, (i) diz que a ausência de resultados é um evento, (ii) afirma que
os resultados possíveis que não estão num determinado evento compõem um evento,
e (iii) afirma que a reunião de uma lista -- possivelmente infinita -- de eventos
forma um evento.

\begin{prop:demorgan}
Se $\mathcal{B}$ é uma $\sigma$-Álgebra do espaço amostral $\Omega$, e
$\{A_1, A_2, \ldots\} \subset \mathcal{B}$, então
\[
	\bigcap_{i = 1}^\infty A_i \in \mathcal{B}
\]
\begin{proof}
	\[
		\bigcup_{i = 1}^\infty A_i^c \in \mathcal{B} \Longleftrightarrow \left(\bigcap_{i=1}^\infty A_i\right)^c
		\in \mathcal{B}
	\]
	o que leva a
	\[
		\bigcap_{i=1}^\infty A_i \in \mathcal{B}
	\]
\end{proof}
\end{prop:demorgan}

\subsection{Axiomas de Kolmogorov}

Para cada evento de um espaço amostral podemos associar uma probabilidade, que consiste
numa medida do evento.
Iremos requerer que essa medida seja não negativa e no máximo 1, de modo a se adequar
a nossa intuição de probabilidade como um número entre 0 e 1.
Além disso, se esperamos que o experimento tenha algum evento, então é razoável presumir
que a probabilidade do próprio espaço amostral é igual a 1.

\begin{def:probabilidade}[Axiomas de Kolmogorov]
Seja $\mathcal{B}$ uma $\sigma$-Álgebra de um conjunto $\Omega$.
Uma medida de probabilidade $P: \mathcal{B} \to \mathbb{R}$ satisfaz:
\begin{enumerate}
	\item[i] $\forall E \in \mathcal{B}, P(E) \geq 0$

	\item $P(\Omega) = 1$

	\item Se $\{E_1, E_2, \ldots\}$ são eventos disjuntos em $\Omega$, então
	      \[
		      P\left(\bigcup_{i=1}^\infty E_i\right) = \sum_{i=1}^\infty P(E_i)
	      \]
\end{enumerate}
\end{def:probabilidade}

Num experimento aleatório, dados os resultados conhecidos, as coleções desses resultados
que compõem eventos e uma medida da probabilidade desses eventos, obtemos, então,
o espaço de probabilidade do experimento.

\begin{def:espaco de probabilidade}
A tripla $(\Omega, \mathcal{B}, P)$ é um espaço de probabilidade se, e somente se,
\begin{enumerate}
	\item[i] $\Omega$ é um espaço amostral
	\item[ii] $\mathcal{B}$ é uma $\sigma$-Álgebra de $\Omega$
	\item[iii] $P: \mathcal{B} \to \mathbb{R}$ é uma medida de probabilidade.
\end{enumerate}
\end{def:espaco de probabilidade}

% \subsection{Eventos Equiprováveis}
% Numa noção intuitiva de que quando o evento não possui vícios, os seus resultados
% são ditos como equiprováveis, como os lados de um dado ou as faces de uma moeda honestos.

% \begin{def:evento equiprovavel}
% Um evento $E \subset \Omega$ é equiprovável se, e somente se, todo subconjunto unitário de
% $E$ possui a mesma probabilidade.
% \end{def:evento equiprovavel}

% \begin{prop:equiprovavel}
% \label{prop:equiprovavel}
% A probabilidade de um evento $E$ num espaço $\Omega$ equiprovável
% é igual a
% \[
% 	\frac{|E|}{|\Omega|}
% \]
% \begin{proof}
% 	Com efeito, temos que
% 	\[
% 		P(E_1) = \ldots = P(E_n)
% 	\]
% 	em que $E_i$ são eventos unitários de $\Omega$, tal que $|\Omega| = n$. Como eles são
% 	disjuntos, então a soma de suas probabilidades deve ser igual a 1, ou seja,
% 	\[
% 		\forall i \in [n],\ \left(P(E_i) = \frac{1}{n}\right).
% 	\]
% 	Segue que
% 	\[
% 		P(E) = |E| \cdot \frac{1}{n} = \frac{|E|}{|\Omega|}
% 	\]
% \end{proof}
% \end{prop:equiprovavel}

% \begin{cor:pie}
% Se $\mathcal{F} = \{E_1, \ldots, E_m\}$ é uma família de eventos quaisquer dum espaço
% $\Omega$ equiprovável, então
% \[
% 	P\left(\bigcup_{E \in \mathcal{F}}E\right) = \sum_{k=1}^m (-1)^{k+1}
% 	\sum_{D \in \Parts{\mathcal{F}}\land |D| = k} P\left(\bigcap_{E \in D} E\right)
% \]
% \begin{proof}
% 	A cardinalidade da união dos conjuntos de $\mathcal{F}$ é dado pelo Princípio da
% 	Exclusão e Inclusão, ou seja
% 	\[
% 		\left|\bigcup_{E \in \mathcal{F}}E\right| = \sum_{k=1}^m (-1)^{k+1}
% 		\sum_{D \in \Parts{\mathcal{F}} \land |D| = k} \left|\bigcap_{E \in D} E\right|
% 	\]
% 	Dividindo a expressão acima por $n = |\Omega|$ obtém-se a probabilidade da união dos
% 	conjuntos de $\mathcal{F}$. Contudo, uma vez que a interseção de eventos é um evento,
% 	então podemos obter a probabilidade da união como
% 	\[
% 		P\left(\bigcup_{E \in \mathcal{F}}E\right) = \sum_{k=1}^m (-1)^{k+1}
% 		\sum_{D \in \Parts{\mathcal{F}}\land |D| = k} \frac{1}{n}\left|\bigcap_{E \in D} E\right|
% 	\]
% 	que equivale a expressão desejada.
% \end{proof}
% \end{cor:pie}

% \subsection{Probabilidade como Função Contínua em um Conjunto}

% \begin{def:eventos encaixados}
% Uma sequência ${E_i}_{i \in \mathbb{N}}$ é encaixado se
% \[
% 	E_1 \subset E_2 \subset \ldots
% \]
% ou
% \[
% 	E_1 \supset E_2 \supset \ldots.
% \]
% O primeiro caso é chamado de crescente, enquanto o segundo de descrescente.
% \end{def:eventos encaixados}

% \begin{prop:eventos encaixados}
% Se ${E_i}_{i \in \mathbb{N}}$ eventos encaixados crescentes, então
% \[
% 	\lim_{n \to \infty} P(E_n) = P\left(\lim_{n \to \infty} E_n\right)
% \]
% \begin{proof}
% 	.
% \end{proof}
% \end{prop:eventos encaixados}
